\section{合同的基本原理}
\begin{enumerate}
  \item 分类要背，是理解各类合同的前提。
  \item 互负对待给付义务。单务，\al{赠予，担保}；\zd{双务}，\al{买卖}。
  \item 法律赋予特定的名称 \hz 设定具体规则。有名/典型，\al{劳动合同}；\zd{无名}，类推适用，\al{培训协议}。
  \item 有没有对价。有偿；\zd{无偿}。\al{赠予、担保、借用；委托、保管、自然人借贷利息} \yjt 可有偿可无偿。
  \item 成立标准。诺成/非要物，一诺既成立；\zd{实践性/要物}，\al{订金、自然人之间的借贷；保管、借用}。
  \item 要不要有特定形式。要式，不要式。
  \item 两个合同的关系（牵连）。预约，本约。\al{意向书，备忘录、认购书}。不得强制要求履行本约。
  \item 射幸，确定。
  \item 一时性，继续性。\al{买卖合同；供电合同、租赁合同}。
\end{enumerate}
